\chapter{Vues Facelets}\label{chap:vues}

\section{Gabarit et composition}
Facelets remplace JSP comme technologie de vue : les pages sont des documents XHTML bien formés, ce qui permet de les
valider avant de lancer le serveur. Le gabarit \code{WEB-INF/templates/layout.xhtml} contient le squelette commun
(en-tête, menu latéral, zone de messages, pied de page) et trois points d'insertion : les métadonnées, le contenu et, au besoin, des
scripts. Chaque page déclare le gabarit avec \code{template=} et remplit ses zones avec \code{ui:define}. Le gabarit est placé sous
\code{WEB-INF} : aucune adresse ne permet de le demander directement.

Les messages JSF s'affichent à un seul endroit, dans le gabarit, avec une boucle sur \code{facesContext.messageList}. La classe
CSS de chaque message dépend de sa gravité.

\lstinputlisting[language=XML,caption={Affichage des messages dans le gabarit},label={lst:messages}]{extraits/layout-messages.txt}

\section{Composants composites}
Cinq fragments répétés dans plusieurs pages sont devenus des composants composites, rangés dans
\code{resources/components} (espace de noms \code{http://xmlns.jcp.org/jsf/composite/components}, préfixe \code{t}).

\begin{table}[htbp]
\centering
\caption{Composants composites}\label{tab:composites}
\begin{tabularx}{\linewidth}{@{}lX@{}}
\toprule
\textbf{Composant} & \textbf{Rôle}\\
\midrule
\code{t:csrf} & Champ caché portant le jeton anti-CSRF ; présent dans chaque formulaire.\\
\code{t:icone} & Icône vectorielle, dessinée à partir d'un tracé fourni par \code{Affichage.icone}.\\
\code{t:visuel} & Image d'un produit, ou pastille d'initiales quand il n'a pas d'image.\\
\code{t:vide} & Bloc « rien à afficher » avec illustration, titre, texte et action facultative.\\
\code{t:champ} & Champ de formulaire (texte, mot de passe ou zone de texte) avec son libellé et l'erreur correspondante, lue dans la table d'erreurs du service.\\
\bottomrule
\end{tabularx}
\end{table}

Un composant composite n'a pas de classe Java ; c'est un modèle paramétré. Sa valeur est de rendre impossible un oubli :
tout formulaire inclut \code{t:csrf}, et un test statique le vérifie (chapitre~\ref{chap:tests}). Comme ces fichiers ne sont
pas des pages, le \code{web.xml} interdit leur accès direct.

Le composant \code{t:champ} traduit un choix de conception : la validation JSF (\code{required}, validateurs) n'est pas
utilisée. Les règles de saisie vivent dans la couche service, comme à l'atelier~2, et renvoient une table
\guillemotleft~champ $\rightarrow$ message~\guillemotright{}. Une seule source de validation évite que la page et le métier divergent.

\section{Composants JSF et politique de sécurité du contenu}
La politique \code{script-src 'self'} (chapitre~\ref{chap:securite}) interdit tout script en ligne. Or plusieurs
composants JSF classiques en génèrent : \code{h:commandLink} produit un lien dont le \code{onclick} appelle
une fonction JavaScript, et \code{f:ajax} injecte du code. Le choix est donc le suivant.

\begin{itemize}
  \item Les actions sont déclenchées par \code{h:commandButton}, un vrai bouton d'envoi qui n'a besoin d'aucun script.
  \item La navigation simple utilise des liens HTML ordinaires, dont l'adresse est calculée par le bean.
  \item La confirmation avant suppression passe par un attribut \code{data-confirm} lu par \code{app.js}, un fichier
        servi par le serveur lui-même.
  \item La recherche est un formulaire \code{GET} en HTML pur, pas un formulaire JSF.
\end{itemize}

Cette contrainte coûte quelques lignes en moins de confort, et apporte une garantie : aucune injection de balise
\code{<script>} ne peut s'exécuter, même si un échappement venait à être oublié. Elle est vérifiée par
\code{VuesChecks}, qui refuse \code{h:commandLink}, \code{f:ajax} et tout attribut \code{on*}.

\section{Pages réalisées}
Seize fichiers \code{.xhtml} composent l'interface. Les figures qui suivent montrent le parcours d'un visiteur, d'un client,
puis d'un administrateur.

\subsection{Vitrine et fiche produit}
La vitrine affiche des cartes de produits, des pastilles de catégories avec leur effectif, un compteur et une
pagination. Selon le profil, le bas de chaque carte propose « Ajouter au panier » (client), « Se connecter pour acheter »
(visiteur), ou rien (administrateur).

\fig{06-vitrine}{Vitrine : cartes de produits et catégories}{fig:vitrine}
\fig{07-vitrine-recherche}{Recherche et filtre conservés dans l'adresse}{fig:recherche}
\fig{08-detail-produit}{Fiche d'un produit}{fig:detail}

\subsection{Connexion et inscription}
En cas d'échec de connexion, un message rouge s'affiche et le champ du mot de passe est vidé. Le formulaire
d'inscription valide tous les champs en une passe et affiche une erreur sous chacun, en plus d'un message général.

\fig{09-connexion-erreur}{Connexion refusée : message d'erreur}{fig:connexion}
\fig{10-inscription-validation}{Inscription : erreurs de validation sous les champs}{fig:inscription}

\subsection{Panier}
Le panier liste les lignes avec leur sous-total, un champ de quantité et le total en dirhams. Lorsqu'une quantité dépasse
le stock, le service lève une exception métier ; le bean la transforme en message et le panier reste inchangé.

\fig{11-panier}{Panier d'un client}{fig:panier}
\fig{12-panier-stock-refuse}{Refus d'une quantité supérieure au stock}{fig:stock}

\subsection{Administration}
Les écrans d'administration sont réservés au rôle \code{ADMIN} : liste des produits avec actions (modifier, retirer de la
vente, supprimer avec confirmation), formulaire produit, gestion des catégories.

\fig{14-admin-produits}{Administration : liste des produits}{fig:admin}
\fig{15-admin-produit-form}{Formulaire produit : erreurs de validation}{fig:adminform}
